\section{再生核希尔伯特空间：把函数当成点}
\label{sec:13-Machine-Learning/09-Kernel-Methods-and-Gaussian-Processes/03}

你把一个 RBF 核的 SVM 交付上线，运维问模型文件里存的是什么、有多大。按上一节的说法，这个模型是在无穷维特征空间里学到的一张超平面，权重向量有无穷多个分量——可导出的文件只有几百 KB，里面是若干个训练样本和同样多的系数。运维追问：无穷维的解为什么用有限个数就写得下？这不是实现上的偷工，而是一条定理的结论。

要把这条定理说清楚，得先回答上一节挂起的两个问题：\(k\) 背后的特征空间到底是什么，正则项 \(\lambda\norm{\cdot}^2\) 惩罚的是哪个范数。本节把核的家——再生核希尔伯特空间（RKHS）——造出来，再证明表示定理。结论会说明：训练损失看不见训练点张成空间之外的分量，正则项却要为它付费，所以最优解只能由训练点上的核函数展开。

\subsection{核函数自己就能张成一个空间}

给定输入集合 \(\mathcal{X}\) 与正定核 \(k\)（任意有限点集的 Gram 矩阵半正定，见上一节的判据）。不借助任何外来的特征映射，直接用 \(k\) 造空间：对每个 \(x\in\mathcal{X}\)，把

\[
k(x,\cdot):\ t\mapsto k(x,t)
\]

看作一个函数，收集它们的全部有限线性组合

\[
\mathcal{H}_0
=\Bigl\{\,f(\cdot)=\sum_{i=1}^{m}\alpha_i\,k(x_i,\cdot):
\ m\in\N,\ \alpha_i\in\R,\ x_i\in\mathcal{X}\,\Bigr\}.
\]

\(\mathcal{H}_0\) 对加法与数乘封闭，是一个向量空间。它的元素是``以某些点为锚的核的叠加''——RBF 核下就是一堆以 \(x_i\) 为中心、高度与符号由 \(\alpha_i\) 决定的钟形曲线之和。下一步是给它内积。

\subsection{一个让核再生的内积}

对 \(f=\sum_{i=1}^{m}\alpha_i k(x_i,\cdot)\) 与 \(g=\sum_{j=1}^{m'}\beta_j k(z_j,\cdot)\)，定义

\[
\ip{f}{g}_{\mathcal{H}_0}
=\sum_{i=1}^{m}\sum_{j=1}^{m'}\alpha_i\beta_j\,k(x_i,z_j).
\]

取 \(g=k(x,\cdot)\)，即 \(m'=1\)、\(\beta_1=1\)、\(z_1=x\)，立刻得到

\[
\ip{f}{k(x,\cdot)}_{\mathcal{H}_0}
=\sum_{i=1}^{m}\alpha_i\,k(x_i,x)
=f(x).
\]

这就是再生性质：与 \(k(x,\cdot)\) 做内积等于在 \(x\) 处求值。函数 \(k(x,\cdot)\) 因此被称为 \(x\) 处的求值代表，它把``求值''这个操作变成了空间里的一个向量。

定义是否合法要检查两点。其一，\(\ip{f}{g}\) 会不会依赖 \(f,g\) 的具体表示？同一个函数可能有多种锚点取法。把定义改写为

\[
\ip{f}{g}_{\mathcal{H}_0}
=\sum_{i=1}^{m}\alpha_i\,g(x_i)
=\sum_{j=1}^{m'}\beta_j\,f(z_j),
\]

第一个等号说明内积的值只依赖 \(g\) 这个函数本身而非它的表示，第二个等号说明只依赖 \(f\) 本身，故内积良定义。其二，它是否真是内积？双线性与对称由定义直接成立；半正定来自核的正定性：对 \(f=\sum_i\alpha_i k(x_i,\cdot)\)，

\[
\norm{f}_{\mathcal{H}_0}^{2}
=\sum_{i=1}^{m}\sum_{j=1}^{m}\alpha_i\alpha_j\,k(x_i,x_j)
=\alpha^{\trans}K\alpha\ge0,
\]

其中 \(K\) 是点集 \(\set{x_i}\) 的 Gram 矩阵。还需排除``范数为零的非零函数''：由半内积的 Cauchy--Schwarz 不等式与再生性质，

\[
\abs{f(x)}
=\abs{\ip{f}{k(x,\cdot)}}
\le\norm{f}_{\mathcal{H}_0}\,\sqrt{k(x,x)},
\]

故 \(\norm{f}=0\) 蕴含 \(f\) 处处为零。内积的全部要求满足。上一节留下的充分性到这里补齐了：半正定的 \(k\) 确实造得出一个内积空间，特征映射可以取 \(\phi(x)=k(x,\cdot)\)。

\subsection{完备化之后，求值仍然连续}

\(\mathcal{H}_0\) 一般不完备：内积下的 Cauchy 列可能收敛不到任何有限线性组合。标准做法是完备化——把所有 Cauchy 列的极限添进来，得到的 Hilbert 空间记为 \(\mathcal{H}_k\)，即 \(k\) 的再生核希尔伯特空间。完备化保持内积与再生性质：对任意 \(f\in\mathcal{H}_k\) 仍有

\[
f(x)=\ip{f}{k(x,\cdot)}_{\mathcal{H}_k}.
\]

严格地说，\(\mathcal{H}_k\) 中的元素最初是 Cauchy 列的等价类，但再生性质使每个等价类对应一个确定的逐点函数：若 \(\norm{f_m-f}\to0\)，由上面的 Cauchy--Schwarz 估计有 \(f_m(x)\to f(x)\) 对每个 \(x\) 成立；在 \(k(x,x)\) 有界的集合上，收敛还是一致的。范数收敛蕴涵逐点收敛，这不是所有函数空间都有的性质。

对比最能说明问题。在 \(L^2\) 中，``函数在一点处的值''根本没有意义：改变一个点上的值不改变 \(L^2\) 中的元素，求值泛函 \(f\mapsto f(x)\) 无定义，更谈不上连续，一个 \(L^2\) 收敛的函数列可以在每个点上都震荡不收敛。RKHS 中求值泛函不仅良定义，还有界：

\[
\abs{f(x)}\le\sqrt{k(x,x)}\,\norm{f}_{\mathcal{H}_k},
\]

其算子范数恰为 \(\sqrt{k(x,x)}\)。``把每个函数当成点、把求值当成内积''的含义全在这里：几何（内积、范数、正交）与分析（求值、逐点收敛）被焊在了一起。

\subsection{范数惩罚的是核认定的不自然变化}

对有限张成的 \(f=\sum_i\alpha_i k(x_i,\cdot)\)，范数已由 Gram 矩阵给出：\(\norm{f}^{2}=\alpha^{\trans}K\alpha\)。但不能只看某组系数的大小来判断函数复杂度，Gram 矩阵奇异时同一个函数甚至可能有多组系数表示。更稳妥的理解来自核积分算子的谱分解：在核的特征函数方向上，低特征值方向要实现同样幅度，通常需要付出更大的 RKHS 范数。核范数因此是``以 \(k\) 的几何为度量的复杂度''。对 RBF 这类平滑核，它确实强烈抑制快速起伏，但不同核的范数不能一概等同于普通导数惩罚。

一天只在早、中、晚看三次温度计，再补出整天的温度曲线时，``曲线应该多平滑''已经是一项假设。长度尺度取得很大，会把相隔数小时的读数也视为强相关，开窗后短暂降温或做饭时的短暂升温可能被抹平；长度尺度很小，又可能把温度计的一次抖动画成真实起伏。核范数惩罚的是所选核认为不自然的变化尺度，脱离核的选择并不存在唯一正确的粗糙度。

空间污染浓度插值把同一选择推向更高风险。RBF 的长度尺度较大（即 \(\gamma\) 较小）时，公路旁很窄的污染峰在 RKHS 范数上罚得很重，容易被平滑掉；长度尺度很小时，模型又可能把传感器噪声解释成真实起伏。\(\lambda\) 只是在\emph{已经选定的核几何内}限制范数，替代不了尺度选择。若决策关心是否越过局部安全阈值，仅用全区域平均预测误差来选核与 \(\lambda\)，可能恰好掩盖最需要保留的那个窄峰。

\subsection{表示定理：无穷维塌缩成 n 维}

现在可以回答运维的问题。设训练数据 \((x_1,y_1),\dots,(x_n,y_n)\)，\(\lambda>0\)，经验损失 \(L\) 任意（只要求问题有极小值），考虑

\[
\min_{f\in\mathcal{H}_k}\
L\bigl(f(x_1),\dots,f(x_n)\bigr)
+\lambda\,\norm{f}_{\mathcal{H}_k}^{2}.
\]

这是一个无穷维优化问题，变量是整个函数。表示定理断言：任何极小解都可写成

\[
f^{*}(\cdot)=\sum_{i=1}^{n}\alpha_i\,k(x_i,\cdot),
\]

即只由 \(n\) 个训练点的锚张成，自由度只有 \(n\) 个系数。

证明只需要正交分解。记 \(S=\operatorname{span}\set{k(x_1,\cdot),\dots,k(x_n,\cdot)}\)，它是有限维子空间，因此任何 \(f\in\mathcal{H}_k\) 唯一分解为 \(f=f_{\parallel}+f_{\perp}\)，其中 \(f_{\parallel}\in S\)，\(f_{\perp}\perp S\)。对每个训练点，由再生性质与正交性，

\[
f(x_i)
=\ip{f}{k(x_i,\cdot)}
=\ip{f_{\parallel}}{k(x_i,\cdot)}+\ip{f_{\perp}}{k(x_i,\cdot)}
=f_{\parallel}(x_i)+0,
\]

因为 \(k(x_i,\cdot)\in S\) 而 \(f_{\perp}\) 与 \(S\) 正交。于是经验损失只看见 \(f_{\parallel}\)：

\[
L\bigl(f(x_1),\dots,f(x_n)\bigr)
=L\bigl(f_{\parallel}(x_1),\dots,f_{\parallel}(x_n)\bigr).
\]

而范数部分满足

\[
\norm{f}^{2}
=\norm{f_{\parallel}}^{2}+\norm{f_{\perp}}^{2}
\ge\norm{f_{\parallel}}^{2},
\]

且 \(f_{\perp}\ne0\) 时严格大于。对任何候选 \(f\)，扔掉垂直分量，损失不变而惩罚严格变小；极小解不可能携带垂直分量，必落在 \(S\) 中。条件对账：证明用到的只有 Hilbert 空间的正交投影、再生性质，以及惩罚项对范数严格单调递增这三件事，损失 \(L\) 的形状从未进入论证。

下图是这三行代数的全部几何内容。

\begin{center}
\begin{tikzpicture}[>=stealth,font=\small]

\draw[black!30] (0,0) ellipse (5.1cm and 2.9cm);
\node[black!45] at (-3.9,2.2) {\(\mathcal{H}_k\)};

\fill[black!8] (-3.3,-1.5) -- (1.5,-1.5) -- (3.3,0.1) -- (-1.5,0.1) -- cycle;
\draw[black!45] (-3.3,-1.5) -- (1.5,-1.5) -- (3.3,0.1) -- (-1.5,0.1) -- cycle;
\node[black!55,font=\footnotesize,anchor=west] at (-3.15,-1.95)
  {\(S=\operatorname{span}\set{k(x_i,\cdot)}\)};

\coordinate (O) at (-0.9,-1.0);
\coordinate (P) at (1.5,-0.55);
\coordinate (F) at (1.5,1.35);

\fill (O) circle (1.6pt);
\node[anchor=north,font=\footnotesize] at (O) {\(0\)};

\draw[->,very thick] (O) -- (P) node[midway,below,font=\footnotesize] {\(f_{\parallel}\)};
\draw[->,very thick] (P) -- (F) node[midway,right,font=\footnotesize] {\(f_{\perp}\)};
\draw[->,dashed] (O) -- (F) node[midway,above left,font=\footnotesize] {\(f\)};
\draw (1.5,-0.55) -- (1.28,-0.55) -- (1.28,-0.33) -- (1.5,-0.33);

\fill (-2.2,-1.05) circle (1.6pt) node[anchor=north,font=\footnotesize] {\(k(x_1,\cdot)\)};
\fill (0.4,-1.15) circle (1.6pt) node[anchor=north,font=\footnotesize] {\(k(x_2,\cdot)\)};

\node[font=\footnotesize,align=left,anchor=west] at (2.0,1.35)
  {损失看不见它\\范数却要付费};

\end{tikzpicture}
\end{center}

\emph{垂直分量对训练损失毫无贡献，对惩罚项却分文不少，于是最优解把它整个丢掉。}

结论值得拆开咀嚼。第一，图中的论证与损失的形状无关——平方损失、hinge 损失、logistic 损失全部适用，这解释了为什么 SVM、核岭回归、核 logistic 回归共享同一个``由 \(n\) 个样本展开''的形态。第二，这是精确降维：无穷维问题的解恰好等于某个 \(n\) 维问题的解，并非近似，剩下的工作只是把目标函数用 Gram 矩阵写出来。第三，它与上一节 Ridge 回归 Gram 形式中 \(w=X^{\trans}\alpha\) 的观察是同一个现象：那里的张成空间由训练样本自身给出，这里由 \(k(x_i,\cdot)\) 给出，普通 Ridge 就是线性核的特例。

\subsection{家已经建好，账单还没结}

回到运维的问题：模型文件里存的是 \(n\) 个训练点与 \(n\) 个系数，因为表示定理保证再多的自由度也用不上。顺带回答的还有本节开头的两问——\(k\) 背后的空间是 \(\mathcal{H}_k\)，其中的函数是核的叠加及其极限；正则项惩罚的是这个空间几何意义下的不光滑程度。

要付的开销则第一次浮出水面：解的规模跟着样本数走，而不是跟着特征维数走。\hyperref[sec:13-Machine-Learning/09-Kernel-Methods-and-Gaussian-Processes/04]{``核岭回归与核方法的代价''}把定理用于平方损失推出闭式解，并把这笔账算清楚。
